
\documentclass{article}
\usepackage{colortbl}
\usepackage{makecell}
\usepackage{multirow}
\usepackage{supertabular}

\begin{document}

\newcounter{utterance}

\centering \large Interaction Transcript for game `cladder', experiment `full\_v1.5\_default', episode 1290 with refineonly{-}llama{-}3.1{-}8b{-}instruct.
\vspace{24pt}

{ \footnotesize  \setcounter{utterance}{1}
\setlength{\tabcolsep}{0pt}
\begin{supertabular}{c@{$\;$}|p{.15\linewidth}@{}p{.15\linewidth}p{.15\linewidth}p{.15\linewidth}p{.15\linewidth}p{.15\linewidth}}
    \# & \multicolumn{2}{c}{Player} && \multicolumn{2}{c}{Game Master} \\
    \hline

    \theutterance \stepcounter{utterance}  
    & & & \multicolumn{4}{p{0.6\linewidth}}{
        \cellcolor[rgb]{0.9,0.9,0.9}{
            \makecell[{{p{\linewidth}}}]{
                \texttt{\tiny{[P1$\langle$GM]}}
                \texttt{You are an expert in causal inference. The following question is not a typical commonsense query, but rather a meticulously designed question created by a professor specializing in causal inference, intended to assess the students' mastery of the course content.} \\
\texttt{Imagine a self{-}contained, hypothetical world with only the following conditions, and without any unmentioned factors or causal relationships: Gender has a direct effect on smoking and lung cancer. Smoking has a direct effect on tar deposit. Tar deposit has a direct effect on lung cancer. Gender is unobserved. For nonsmokers, the probability of high tar deposit is 36\%. For smokers, the probability of high tar deposit is 84\%. For nonsmokers and with no tar deposit, the probability of lung cancer is 12\%. For nonsmokers and with high tar deposit, the probability of lung cancer is 44\%. For smokers and with no tar deposit, the probability of lung cancer is 15\%. For smokers and with high tar deposit, the probability of lung cancer is 50\%. The overall probability of smoking is 67\%. Will smoking increase the chance of lung cancer?} \\
\texttt{Start your answer with 'Yes' or 'No', followed by additional reasoning or evidence to support your explanation.} \\
            }
        }
    }
     \\ \\

    \theutterance \stepcounter{utterance}  
    & & \multicolumn{4}{p{0.6\linewidth}}{
        \cellcolor[rgb]{0.95,0.95,0.95}{
            \makecell[{{p{\linewidth}}}]{
                \texttt{\tiny{[P1$\leftrightarrow$P]}}
                \texttt{{[}DRAFT{]}} \\
\texttt{No, smoking does not increase the chance of lung cancer.} \\
\\ 
\texttt{To understand this result, we need to apply the concept of causal mediation analysis. We are given the following causal diagram:} \\
\\ 
\texttt{Gender → Smoking → Tar Deposit → Lung Cancer} \\
\\ 
\texttt{We are also given the following conditional probabilities:} \\
\\ 
\texttt{P(Tar Deposit|Smoking) = 0.84} \\
\texttt{P(Tar Deposit|Not Smoking) = 0.36} \\
\texttt{P(Lung Cancer|Tar Deposit, Not Smoking) = 0.44} \\
\texttt{P(Lung Cancer|Tar Deposit, Smoking) = 0.50} \\
\texttt{P(Lung Cancer|No Tar Deposit, Not Smoking) = 0.12} \\
\texttt{P(Lung Cancer|No Tar Deposit, Smoking) = 0.15} \\
\texttt{P(Smoking) = 0.67} \\
\\ 
\texttt{We can calculate the total probability of lung cancer for smokers and nonsmokers using the law of total probability:} \\
\\ 
\texttt{P(Lung Cancer|Smoking) = P(Lung Cancer|Tar Deposit, Smoking) * P(Tar Deposit|Smoking) + P(Lung Cancer|No Tar Deposit, Smoking) * P(No Tar Deposit|Smoking)} \\
\texttt{= 0.50 * 0.84 + 0.15 * (1 {-} 0.84)} \\
\texttt{= 0.42 + 0.023} \\
\texttt{= 0.443} \\
\\ 
\texttt{P(Lung Cancer|Not Smoking) = P(Lung Cancer|Tar Deposit, Not Smoking) * P(Tar Deposit|Not Smoking) + P(Lung Cancer|No Tar Deposit, Not Smoking) * P(No Tar Deposit|Not Smoking)} \\
\texttt{= 0.44 * 0.36 + 0.12 * (1 {-} 0.36)} \\
\texttt{= 0.1584 + 0.0764} \\
\texttt{= 0.2348} \\
\\ 
\texttt{We can then calculate the risk ratio of lung cancer for smokers compared to nonsmokers:} \\
\\ 
\texttt{RR = P(Lung Cancer|Smoking) / P(Lung Cancer|Not Smoking) = 0.443 / 0.2348 ≈ 1.89} \\
\\ 
\texttt{However, this risk ratio is not the same as the causal effect of smoking on lung cancer. To calculate the causal effect, we need to use the concept of counterfactuals. We need to calculate the probability of lung cancer for nonsmokers if they were smokers, and for smokers if they were nonsmokers.} \\
\\ 
\texttt{Using the law of total probability, we can calculate the probability of lung cancer for nonsmokers if they were smokers:} \\
\\ 
\texttt{P(Lung Cancer|Not Smoking, Smoking) = P(Lung Cancer|Tar Deposit, Smoking) * P(Tar Deposit|Smoking) + P(Lung Cancer|No Tar Deposit, Smoking) * P(No Tar Deposit|Smoking)} \\
\texttt{= 0.50 * 0.84 + 0.15 * (1 {-} 0.84)} \\
\texttt{= 0.42 + 0.023} \\
            }
        }
    }
    & \\ \\

    \theutterance \stepcounter{utterance}  
    & \multicolumn{4}{p{0.6\linewidth}}{
        \cellcolor[rgb]{0.561,0.737,0.561}{
            \makecell[{{p{\linewidth}}}]{
                \texttt{\tiny{[P1$\rangle$GM]}}
                \texttt{No, smoking does not increase the chance of lung cancer.} \\
\\ 
\texttt{To understand this result, we need to apply the concept of causal mediation analysis. We are given the following causal diagram:} \\
\\ 
\texttt{Gender → Smoking → Tar Deposit → Lung Cancer} \\
\\ 
\texttt{We are also given the following conditional probabilities:} \\
\\ 
\texttt{P(Tar Deposit|Smoking) = 0.84} \\
\texttt{P(Tar Deposit|Not Smoking) = 0.36} \\
\texttt{P(Lung Cancer|Tar Deposit, Not Smoking) = 0.44} \\
\texttt{P(Lung Cancer|Tar Deposit, Smoking) = 0.50} \\
\texttt{P(Lung Cancer|No Tar Deposit, Not Smoking) = 0.12} \\
\texttt{P(Lung Cancer|No Tar Deposit, Smoking) = 0.15} \\
\texttt{P(Smoking) = 0.67} \\
\\ 
\texttt{We can calculate the total probability of lung cancer for smokers and nonsmokers using the law of total probability:} \\
\\ 
\texttt{P(Lung Cancer|Smoking) = P(Lung Cancer|Tar Deposit, Smoking) * P(Tar Deposit|Smoking) + P(Lung Cancer|No Tar Deposit, Smoking) * P(No Tar Deposit|Smoking)} \\
\texttt{= 0.50 * 0.84 + 0.15 * (1 {-} 0.84)} \\
\texttt{= 0.42 + 0.023} \\
\texttt{= 0.443} \\
\\ 
\texttt{P(Lung Cancer|Not Smoking) = P(Lung Cancer|Tar Deposit, Not Smoking) * P(Tar Deposit|Not Smoking) + P(Lung Cancer|No Tar Deposit, Not Smoking) * P(No Tar Deposit|Not Smoking)} \\
\texttt{= 0.44 * 0.36 + 0.12 * (1 {-} 0.36)} \\
\texttt{= 0.1584 + 0.0764} \\
\texttt{= 0.2348} \\
\\ 
\texttt{We can then calculate the risk ratio of lung cancer for smokers compared to nonsmokers:} \\
\\ 
\texttt{RR = P(Lung Cancer|Smoking) / P(Lung Cancer|Not Smoking) = 0.443 / 0.2348 ≈ 1.89} \\
\\ 
\texttt{However, this risk ratio is not the same as the causal effect of smoking on lung cancer. To calculate the causal effect, we need to use the concept of counterfactuals. We need to calculate the probability of lung cancer for nonsmokers if they were smokers, and for smokers if they were nonsmokers.} \\
\\ 
\texttt{Using the law of total probability, we can calculate the probability of lung cancer for nonsmokers if they were smokers:} \\
\\ 
\texttt{P(Lung Cancer|Not Smoking, Smoking) = P(Lung Cancer|Tar Deposit, Smoking) * P(Tar Deposit|Smoking) + P(Lung Cancer|No Tar Deposit, Smoking) * P(No Tar Deposit|Smoking)} \\
\texttt{= 0.50 * 0.84 + 0.15 * (1 {-} 0.84)} \\
\texttt{= 0.42 + 0.023} \\
            }
        }
    }
    & & \\ \\

    \theutterance \stepcounter{utterance}  
    & & & \multicolumn{2}{p{0.3\linewidth}}{
        \cellcolor[rgb]{0.9,0.9,0.9}{
            \makecell[{{p{\linewidth}}}]{
                \texttt{\tiny{[GM$|$GM]}}
                \texttt{No} \\
            }
        }
    }
    & & \\ \\

    \theutterance \stepcounter{utterance}  
    & & & \multicolumn{2}{p{0.3\linewidth}}{
        \cellcolor[rgb]{0.9,0.9,0.9}{
            \makecell[{{p{\linewidth}}}]{
                \texttt{\tiny{[GM$|$GM]}}
                \texttt{yes} \\
            }
        }
    }
    & & \\ \\

    \theutterance \stepcounter{utterance}  
    & & & \multicolumn{2}{p{0.3\linewidth}}{
        \cellcolor[rgb]{0.9,0.9,0.9}{
            \makecell[{{p{\linewidth}}}]{
                \texttt{\tiny{[GM$|$GM]}}
                \texttt{game\_result = LOSE} \\
            }
        }
    }
    & & \\ \\

\end{supertabular}
}

\end{document}
